\section{Parameter-Free Second-Order Methods}
\label{sec:parameter-free}

Algorithm~\ref{alg:ncsc-utr} achieves the desired
$\epsilon^{-3/2}$ leading outer dependence and a sharper total
inner-gradient bound, but its outer scale $M_P$ and residual target
$\tau_{\rm out}$ depend on the unknown regularity constants
$\ell$, $\mu$, and $\rho$. We now remove these quantities from the
executable method.

We give two parameter-free constructions. The first combines an
adaptive UTR scheme with uniformly accurate inner solves. It retains
the $\epsilon^{-3/2}$ leading outer dependence, but its total
inner-gradient complexity contains a multiplicative
$\log(1/\epsilon)$. The second construction instead tracks the
sequence of inner maximizers through persistent SCAR calls and adapts
the inner accuracy to the current outer scale. This removes the
target-accuracy logarithm from the leading inner-gradient term.

Both results below are class-uniform asymptotic statements. After
fixing common constants $(\ell,\mu,\rho)$, there exists an unknown
threshold $\epsilon_0(\ell,\mu,\rho)>0$ such that the conclusions hold
simultaneously for every payoff satisfying
Assumption~\ref{ass:ncsc} with these constants and every
initialization. The algorithms themselves do not use these constants
and cannot determine whether a requested tolerance lies below the
corresponding threshold.


\subsection{A Uniform-Accuracy Parameter-Free Method}
\label{sec:uniform-parameter-free}

We first recall the parameter-free strongly convex solver used for
the inner maximization. For fixed $x$, define
\[
    q_x(y):=-f(x,y).
\]
Under Assumption~\ref{ass:ncsc}, $q_x$ is $\mu$-strongly convex and
has an $\ell$-Lipschitz gradient.


\paragraph{SCAR as an inner residual solver.}

We use the following result of \citet{lan2026optimal}.

\begin{theorem}[SCAR; \citet{lan2026optimal}]
\label{thm:scar-blackbox}
Let $\phi:\mathbb{R}^m\to\mathbb{R}$ be differentiable,
$\mu$-strongly convex, and have an $L$-Lipschitz gradient, with
$0<\mu<L$. Let $\tau>0$ and $u_0\in\mathbb{R}^m$.
If $\|\nabla\phi(u_0)\|\le\tau$, return $u_0$.
Otherwise, initialize SCAR with
\[
    \mu_0\ge\mu,
    \qquad
    0<M_0\le4L.
\]
Then SCAR returns $\widehat u$ satisfying
\begin{equation}
    \|\nabla\phi(\widehat u)\|\le\tau
    \label{eq:scar-residual}
\end{equation}
using at most
\begin{equation}
    (4+8\sqrt5\,C_1)
    \sqrt{\frac{L}{\mu}}
    \left[
        \left\lceil
            \log_4\frac{\mu_0}{\mu}
        \right\rceil
        +
        \left\lceil
            \log_2
            \frac{\|\nabla\phi(u_0)\|}{\tau}
        \right\rceil
    \right]
    \label{eq:scar-original-complexity}
\end{equation}
gradient evaluations, where $C_1$ is the universal constant in
Theorem~4.1 of \citet{lan2026optimal}. The line-search estimate
returned by SCAR is at most $2L$.
\end{theorem}

The cited theorem is stated for the strict case $\mu<L$.
The endpoint $L=\mu$ can be handled separately by the same AR
estimate and contributes only a numerical constant; see
Appendix~\ref{app:uniform-parameter-free-proofs}.

SCAR can be initialized without knowing either curvature parameter.
Given $x$, an initial point $y^{\rm in}$, and any
$z\ne y^{\rm in}$, define the observable secant
\begin{equation}
    a_x
    :=
    \frac{
        \|\nabla_yf(x,z)-\nabla_yf(x,y^{\rm in})\|
    }{
        \|z-y^{\rm in}\|
    }.
    \label{eq:scar-secant}
\end{equation}

\begin{corollary}[Parameter-free SCAR residual module]
\label{cor:scar-inner}
Under Assumption~\ref{ass:ncsc}, run SCAR on
$q_x=-f(x,\cdot)$ with
\[
    \mu_0=M_0=a_x
\]
and residual tolerance $\tau$.
No value of $\ell$, $\mu$, or $\kappa$ is supplied to the inner
solver, and the returned point $\widehat y$ satisfies
\[
    \|\nabla_yf(x,\widehat y)\|\le\tau.
\]
Including the secant initialization, the number of inner-gradient
evaluations is
\begin{equation}
    O\!\left(
        \sqrt\kappa
        \left[
            1+\log\kappa
            +
            \log_+
            \frac{
                \|\nabla_yf(x,y^{\rm in})\|
            }{
                \tau
            }
        \right]
    \right).
    \label{eq:scar-inner-complexity}
\end{equation}
Hence secant-initialized SCAR realizes the residual interface in
Definition~\ref{def:inner-residual-module} without using the problem
regularity constants.
\end{corollary}

The proof is deferred to
Appendix~\ref{app:uniform-parameter-free-proofs}.


\paragraph{Uniform inner accuracy.}

The first parameter-free construction uses the same absolute residual
target at every inner call:
\begin{equation}
    \tau_{\rm pf}:=\epsilon^2,
    \qquad
    0<\epsilon<1.
    \label{eq:pf-uniform-residual}
\end{equation}
The nonmonotone value allowance used below is $\epsilon^3$.

Although these quantities do not involve unknown regularity
parameters, sufficiently small $\epsilon$ is still required in the
analysis to convert the observable residual into the desired
primal-function oracle accuracy.

\begin{lemma}[Class-uniform asymptotic oracle accuracy]
\label{lem:pf-asymptotic-oracle}
Fix $\ell\ge\mu>0$ and $\rho>0$, and define
\[
    \kappa:=\frac{\ell}{\mu},
    \qquad
    M_P:=4\sqrt2\,\kappa^3\rho.
\]
There exists
$\epsilon_{\rm or}=\epsilon_{\rm or}(\ell,\mu,\rho)>0$
such that, simultaneously for every payoff satisfying
Assumption~\ref{ass:ncsc} with these common constants and every
$0<\epsilon\le\epsilon_{\rm or}$, any point satisfying
\[
    \|\nabla_yf(x,y)\|\le\epsilon^2
\]
also satisfies
\begin{subequations}
\label{eq:pf-oracle-bounds}
\begin{align}
    \|\widehat g(x,y)-\nabla P(x)\|
    &\le
    \frac{\epsilon}{256},
    \label{eq:pf-gradient-error}
    \\
    \|\widehat H(x,y)-\nabla^2P(x)\|
    &\le
    \frac1{64}
    \sqrt{\frac{M_P\epsilon}{6}},
    \label{eq:pf-hessian-error}
    \\
    |\widehat P(x,y)-P(x)|
    &\le
    \frac{\epsilon^3}{8}.
    \label{eq:pf-value-error}
\end{align}
\end{subequations}
\end{lemma}

The threshold is used only in the analysis and is not supplied to the
algorithm.


\paragraph{Algorithm 2.}

At a current stored tuple $(x,y)$ satisfying
\eqref{eq:pf-uniform-residual}, let
\[
    (\widehat P,g,H)
    :=
    \bigl(
        \widehat P(x,y),
        \widehat g(x,y),
        \widehat H(x,y)
    \bigr),
    \qquad
    h:=\max\{\|g\|,\epsilon\}.
\]
For a trial scale $\sigma\ge\epsilon$, consider
\begin{equation}
\begin{aligned}
    \min_{\|d\|\le\sqrt h/(4\sigma)}
    \quad
    g^\top d
    +
    \frac12d^\top
    \left(
        H
        +
        \frac{\sigma\sqrt\epsilon}{64}I
        +
        \sigma\sqrt h\,I
    \right)d.
\end{aligned}
\label{eq:pf-robust-subproblem}
\end{equation}

\begin{algorithm}[t]
\DontPrintSemicolon
\caption{Uniform-accuracy parameter-free NC--SC UTR}
\label{alg:pf-residual-utr}
\KwIn{
$x_0,y_{-1}$, $0<\epsilon<1$, and a parameter-free inner residual
module $\mathcal S$ satisfying \eqref{eq:maxres-interface}
}
$x\leftarrow x_0$,
$\sigma\leftarrow\epsilon$,
$\eta\leftarrow1/256$\;

$y\leftarrow
\operatorname{MaxRes}_{\mathcal S}
(x_0,y_{-1},\epsilon^2)$\;

compute
$(\widehat P,g,H)$ from
\eqref{eq:corrected-value} and
\eqref{eq:corrected-derivatives}\;

\While{true}{
    $h\leftarrow\max\{\|g\|,\epsilon\}$\;

    \If{
        $\|g\|\le\epsilon$ and
        $\lambda_{\min}(H)
        \ge-(9/8)\sigma\sqrt\epsilon$
    }{
        \Return{$x$}\;
    }

    solve \eqref{eq:pf-robust-subproblem} globally and obtain
    $(d,\lambda)$;
    $r\leftarrow\|d\|$\;

    \If{$r=0$}{
        \Return{$x$}\;
    }

    $x^+\leftarrow x+d$\;

    $y^+\leftarrow
    \operatorname{MaxRes}_{\mathcal S}
    (x^+,y,\epsilon^2)$\;

    compute
    $(\widehat P^+,g^+,H^+)$ at $(x^+,y^+)$\;

    \eIf{
        $r<\sqrt h/(4\sigma)$,
        $\widehat P^+\le\widehat P+\epsilon^3$, and
        $\|g^+\|\le h/2$
    }{
        accept the trial\;
    }{
        \eIf{
            $r=\sqrt h/(4\sigma)$,
            $\widehat P-\widehat P^+
            \ge
            \eta\bigl(h^{3/2}/\sigma+\lambda r^2\bigr)
            +\epsilon^3$,
            and
            $\|g^+\|\le h/2+\lambda r$
        }{
            accept the trial\;
        }{
            reject the trial\;
        }
    }

    \eIf{the trial is accepted}{
        $(x,y,\widehat P,g,H)
        \leftarrow
        (x^+,y^+,\widehat P^+,g^+,H^+)$\;
        $\sigma\leftarrow\max\{\epsilon,\sigma/2\}$\;
    }{
        discard the candidate tuple\;
        $\sigma\leftarrow2\sigma$\;
    }
}
\end{algorithm}

Every quantity tested by Algorithm~\ref{alg:pf-residual-utr} is
observable. In particular, if $\mathcal S$ is chosen as the
secant-initialized SCAR module in Corollary~\ref{cor:scar-inner},
the executable method uses none of
\[
    \ell,\quad
    \mu,\quad
    \rho,\quad
    \kappa,\quad
    M_P,\quad
    \Delta_P.
\]
The fixed inner target $\epsilon^2$ is independent of the trial scale,
so rejection changes only the outer trust-region model; the corrected
tuple at the current point remains unchanged.

The proof uses the fact that once
\[
    \sigma^2\ge\frac{M_P}{6},
\]
the Hessian padding is large enough for the inexact trust-region
estimate of Lemma~\ref{lem:common-inexact-tr-step}. A finite-prefix
counting argument then controls accepted and rejected trials.
The detailed safe-scale and counting lemmas are deferred to
Appendix~\ref{app:uniform-parameter-free-proofs}.

\begin{theorem}[Uniform-accuracy parameter-free complexity]
\label{thm:pf-scar-utr}
Fix $\ell\ge\mu>0$ and $\rho>0$.
There exists
\[
    \epsilon_0
    =
    \epsilon_0(\ell,\mu,\rho)
    \in(0,1)
\]
such that the following statement holds simultaneously for every
payoff satisfying Assumption~\ref{ass:ncsc} with these common
constants and every initialization $(x_0,y_{-1})$.

Define, only for the purpose of stating the bound,
\[
    \kappa:=\frac{\ell}{\mu},
    \qquad
    M_P:=4\sqrt2\,\kappa^3\rho,
    \qquad
    \Delta_P:=P(x_0)-P^*.
\]
For every $0<\epsilon\le\epsilon_0$,
Algorithm~\ref{alg:pf-residual-utr} terminates and returns
$\widehat x$ satisfying
\begin{equation}
    \|\nabla P(\widehat x)\|
    \le
    \frac{257}{256}\epsilon,
    \qquad
    \lambda_{\min}(\nabla^2P(\widehat x))
    \ge
    -\sqrt{M_P\epsilon}.
    \label{eq:pf-return}
\end{equation}
Consequently, using internal tolerance $256\epsilon/257$ gives the
standard
$(\epsilon,\sqrt{M_P\epsilon})$ guarantee.

If $Q$ denotes the total number of solved robust trust-region
subproblems, then
\begin{equation}
    Q
    =
    O\!\left(
        \Delta_P\sqrt{M_P}\,
        \epsilon^{-3/2}
        +
        \log_+
        \frac{\|\nabla P(x_0)\|}{\epsilon}
        +
        \log_+
        \frac{\sqrt{M_P}}{\epsilon}
        +1
    \right).
    \label{eq:pf-outer-complexity}
\end{equation}
Thus its leading outer dependence is
\begin{equation}
    O\!\left(
        \Delta_P
        \kappa^{3/2}\sqrt\rho\,
        \epsilon^{-3/2}
    \right).
    \label{eq:pf-outer-leading}
\end{equation}
\end{theorem}

This theorem is an asymptotic parameter-free result in the precise
class-uniform sense stated above: the algorithm does not use the
regularity constants, but the theorem does not provide an observable
certificate that a given $\epsilon$ lies below
$\epsilon_0(\ell,\mu,\rho)$.

We next count the work of the concrete SCAR implementation.

\begin{corollary}[Uniform-accuracy SCAR complexity]
\label{cor:pf-scar-inner-complexity}
Under the hypotheses of
Theorem~\ref{thm:pf-scar-utr}, choose
$\mathcal S$ to be the secant-initialized SCAR module in
Corollary~\ref{cor:scar-inner}.
Let $Q$ be the number of robust trust-region subproblems and define
\begin{equation}
    G_Q
    :=
    1+
    \sqrt{
        2(1+\kappa)\ell
        \left(
            \Delta_P+\frac54Q\epsilon^3
        \right)
    }.
    \label{eq:pf-gradient-level}
\end{equation}
Then the total number $N_y$ of inner-gradient evaluations satisfies
\begin{align}
    N_y
    &=
    O\!\left(
        \sqrt\kappa\,(Q+1)
        \left[
            1+\log\kappa
            +
            \log\!\left(
                1+
                \frac{
                    \ell\sqrt{G_Q}
                }{
                    4\epsilon^3
                }
            \right)
        \right]
    \right)
    \notag\\
    &\quad+
    O\!\left(
        \sqrt\kappa
        \left[
            1+\log\kappa
            +
            \log_+
            \frac{
                \|\nabla_yf(x_0,y_{-1})\|
            }{
                \epsilon^2
            }
        \right]
    \right).
    \label{eq:pf-inner-complexity}
\end{align}
In particular, the uniformly accurate SCAR calls introduce a
multiplicative $\log(1/\epsilon)$ in the leading target-accuracy
dependence of the total inner-gradient count.
\end{corollary}

Thus Algorithm~\ref{alg:pf-residual-utr} already removes all unknown
regularity constants from the executable method while retaining the
desired leading outer rate. Its remaining weakness is not outer
adaptivity, but the decision to solve every candidate inner problem
to the same absolute residual $\epsilon^2$.


\subsection{Persistent Inner Tracking}
\label{sec:persistent-inner-tracking}

We now replace independent high-accuracy inner solves by tracking the
sequence of inner maximizers. The central idea is to retain the
curvature calibration learned by SCAR across successive objectives
\[
    q_x(y)=-f(x,y),
\]
while restarting accelerated momentum whenever the outer point
changes.

Rather than solving every inner problem independently to a prescribed
absolute accuracy, the method maintains only the accuracy required by
the current outer scale.


\paragraph{Persistent SCAR.}

We first record the sequential SCAR property used by the algorithm.

\begin{lemma}[Persistent SCAR halving]
\label{lem:scar-persistent-half}
Let $\{\phi_j\}_{j\ge1}$ be any sequence of differentiable functions
that are all $\mu$-strongly convex and have $L$-Lipschitz gradients.
Set $\kappa:=L/\mu$.

Suppose the retained strong-convexity guess $\nu_0$ and line-search
estimate $\mathcal M_0$ satisfy
\[
    \mu\le\nu_0\le L,
    \qquad
    0<\mathcal M_0\le4L.
\]
A persistent halving call at $(\phi_j,u)$ freezes
\[
    R:=\|\nabla\phi_j(u)\|
\]
and repeatedly applies the AR call used by SCAR until it obtains
$v$ satisfying
\[
    \|\nabla\phi_j(v)\|\le\frac R2.
\]
A failed curvature guess replaces $\nu$ by $\nu/4$ while retaining
the updated line-search estimate.

Across any $N$ successful halving calls, the total number of failed
curvature guesses is at most
\begin{equation}
    \left\lceil
        \log_4\frac{\nu_0}{\mu}
    \right\rceil
    \le
    \lceil\log_4\kappa\rceil,
    \label{eq:persistent-scar-failures}
\end{equation}
and the total number of gradient evaluations is at most
\begin{equation}
    (4+2\sqrt{10}\,C_1)
    \sqrt\kappa
    \left(
        N+
        \left\lceil
            \log_4\frac{\nu_0}{\mu}
        \right\rceil
    \right).
    \label{eq:persistent-scar-count}
\end{equation}
After any nontrivial AR call, the retained line-search estimate is at
most $2L$.
\end{lemma}

The crucial distinction from restarted SCAR is that the
$\log\kappa$ curvature-calibration cost in
\eqref{eq:persistent-scar-failures} is paid across the complete
sequence, rather than independently at every outer call.


\paragraph{Working accuracy and tracking.}

Define
\begin{equation}
    \mathcal A_\epsilon
    :=
    \log(e/\epsilon),
    \qquad
    \sigma_0
    :=
    \mathcal A_\epsilon^{-1},
    \qquad
    c_0:=2^{-12},
    \qquad
    r_\epsilon(\sigma)
    :=
    \frac{\sqrt\epsilon}{4\sigma}.
    \label{eq:gs-parameters}
\end{equation}
The quantity $r_\epsilon(\sigma)$ is an inner accuracy scale and
should not be confused with the trust-region step length.

A stored tuple $(x,y)$ is called \emph{working-accurate} at scale
$\sigma$ if
\begin{equation}
    \|\nabla_yf(x,y)\|
    \le
    c_0\ell r_\epsilon(\sigma).
    \label{eq:gs-working-invariant}
\end{equation}
The constant $\ell$ appears only in the analysis; the executable
method maintains this invariant through observable residual halvings.

\begin{lemma}[Variable-step tracking]
\label{lem:gs-tracking}
Suppose \eqref{eq:gs-working-invariant} holds at $(x,y)$, and let
$d$ be an outer step with $r:=\|d\|$.
Starting from $y$ at $x+d$, apply
\begin{equation}
    N_{\rm tr}(r,\sigma)
    :=
    \left\lceil
        \log_2\left(
            1+
            \frac{r}{
                c_0r_\epsilon(\sigma)
            }
        \right)
    \right\rceil
    \label{eq:gs-tracking-count}
\end{equation}
successful persistent SCAR halvings.
The resulting point satisfies
\eqref{eq:gs-working-invariant} at $x+d$ with the same scale.
At an unchanged outer point, one successful halving restores the
invariant after $\sigma$ is doubled.
\end{lemma}

For later use, set
\begin{equation}
    C_g
    :=
    \frac{
        c_0^2(1+\kappa)\rho\kappa^2
    }{32},
    \qquad
    C_v
    :=
    \frac{
        7c_0^3\rho\kappa^3
    }{1536}.
    \label{eq:gs-error-constants}
\end{equation}

\begin{lemma}[Working-oracle accuracy]
\label{lem:gs-working-errors}
Fix common constants $(\ell,\mu,\rho)$.
For all sufficiently small $\epsilon$, with a threshold depending
only on these constants, every working-accurate tuple at
$\sigma\ge\sigma_0$ satisfies
\begin{equation}
    \|\widehat g-\nabla P(x)\|
    \le
    C_g\frac{\epsilon}{\sigma^2},
    \qquad
    |\widehat P-P(x)|
    \le
    C_v\frac{\epsilon^{3/2}}{\sigma^3}.
    \label{eq:gs-working-errors}
\end{equation}
If additionally
\[
    \sigma^2\ge\frac{M_P}{6},
\]
then
\begin{equation}
\begin{aligned}
    \|\widehat g-\nabla P(x)\|
    &\le
    \epsilon/1024,
    \\
    \|\widehat H-\nabla^2P(x)\|
    &\le
    \sigma\sqrt\epsilon/64,
    \\
    |\widehat P-P(x)|
    &\le
    \epsilon^{3/2}/(4096\sigma).
\end{aligned}
\label{eq:working-safe-errors}
\end{equation}

Furthermore, for any secant $a_x$ of the form
\eqref{eq:scar-secant}, any validation point satisfying
\begin{equation}
    \|\nabla_yf(x,y_{\rm v})\|
    \le
    c_0a_x\epsilon^{3/2}
    \label{eq:validator-residual}
\end{equation}
obeys
\begin{equation}
\begin{aligned}
    \|\widehat g(x,y_{\rm v})-\nabla P(x)\|
    &\le
    \epsilon/1024,
    \\
    \|\widehat H(x,y_{\rm v})-\nabla^2P(x)\|
    &\le
    \sigma\sqrt\epsilon/32
\end{aligned}
\label{eq:validator-errors}
\end{equation}
for every $\sigma\ge\sigma_0$.
\end{lemma}

All three preceding lemmas are proved in
Appendix~\ref{app:persistent-scar-proofs}.


\paragraph{Algorithm 3.}

Given a working-accurate tuple, define
\begin{equation}
    h:=\max\{\|\widehat g\|,\epsilon\},
    \qquad
    R(h,\sigma)
    :=
    \frac{\sqrt h}{4\sigma},
    \label{eq:gs-radius}
\end{equation}
and solve
\begin{equation}
    \min_{\|d\|\le R(h,\sigma)}
    \quad
    \widehat g^\top d
    +
    \frac12d^\top
    \left(
        \widehat H
        +
        \frac{\sigma\sqrt\epsilon}{64}I
        +
        \sigma\sqrt h\,I
    \right)d.
    \label{eq:gs-subproblem}
\end{equation}
For a global solution $d$ and its multiplier $\lambda\ge0$, write
\begin{equation}
    r:=\|d\|,
    \qquad
    W:=\sigma\sqrt h\,r^2,
    \qquad
    U:=\lambda r^2,
    \qquad
    T:=\frac{\epsilon^{3/2}}{\sigma}.
    \label{eq:gs-step-quantities}
\end{equation}

\begin{algorithm}[t]
\DontPrintSemicolon
\caption{Parameter-free gradient-scaled TR with persistent SCAR}
\label{alg:gradient-scaled-scar}
\KwIn{$x_0,y_{-1}$ and $0<\epsilon<1$}

$\mathcal A_\epsilon\leftarrow\log(e/\epsilon)$,
$\sigma\leftarrow\mathcal A_\epsilon^{-1}$,
$c_0\leftarrow2^{-12}$,
$x\leftarrow x_0$\;

choose $z_0\ne y_{-1}$ and compute $a_0$ from
\eqref{eq:scar-secant}\;

initialize the persistent SCAR state with
$\nu_0=\mathcal M_0=a_0$\;

at $x_0$, apply persistent SCAR halving from $y_{-1}$ until
\[
    \|\nabla_yf(x_0,y)\|
    \le
    c_0a_0r_\epsilon(\sigma);
\]
compute $(\widehat P,\widehat g,\widehat H)$\;

\While{true}{
    $h\leftarrow\max\{\|\widehat g\|,\epsilon\}$\;

    solve \eqref{eq:gs-subproblem} globally and obtain
    $(d,\lambda)$\;

    $x^+\leftarrow x+d$,
    $r\leftarrow\|d\|$\;

    form $(W,U,T)$ from \eqref{eq:gs-step-quantities}\;

    starting from $y$, apply
    $N_{\rm tr}(r,\sigma)$ successful persistent SCAR halvings
    to $q_{x^+}$ and denote the result by $y^+$\;

    \eIf{$h=\epsilon$ and $r<R(h,\sigma)$}{
        choose $z^+\ne y^+$ and compute the secant $a^+$ at $x^+$\;

        continue persistent SCAR halving at $x^+$ until
        \[
            \|\nabla_yf(x^+,y_{\rm v})\|
            \le
            c_0a^+\epsilon^{3/2};
        \]

        compute
        $(\widehat g_{\rm v},\widehat H_{\rm v})$\;

        \eIf{
            $\|\widehat g_{\rm v}\|\le\epsilon/2$
            and
            $\lambda_{\min}(\widehat H_{\rm v})
            \ge-(21/8)\sigma\sqrt\epsilon$
        }{
            \Return{$x^+$}\;
        }{
            discard the candidate and validator tuples\;
            $\sigma\leftarrow2\sigma$\;
            apply one successful persistent halving at the unchanged $x$\;
            replace $y$ by the returned point and recompute
            $(\widehat P,\widehat g,\widehat H)$\;
        }
    }{
        compute
        $(\widehat P^+,\widehat g^+,\widehat H^+)$
        at $(x^+,y^+)$\;

        \eIf{
            $\widehat P-\widehat P^+
            \ge
            W/8+U/4-T/1024$
            and
            $\|\widehat g^+\|\le h/2+\lambda r$
        }{
            $(x,y,\widehat P,\widehat g,\widehat H)
            \leftarrow
            (x^+,y^+,\widehat P^+,\widehat g^+,\widehat H^+)$\;
        }{
            discard the candidate tuple\;
            $\sigma\leftarrow2\sigma$\;
            apply one successful persistent halving at the unchanged $x$\;
            replace $y$ by the returned point and recompute
            $(\widehat P,\widehat g,\widehat H)$\;
        }
    }
}
\end{algorithm}

The scale in Algorithm~\ref{alg:gradient-scaled-scar} is
nondecreasing. An accepted candidate tuple is stored without further
inner refinement. The persistent SCAR state $(\nu,\mathcal M)$ is
retained after accepted trials, rejected trials, and failed
validations, while accelerated momentum is restarted whenever the
inner objective changes.


\paragraph{Safe scale and accumulated work.}

At every sufficiently small target tolerance, once
\[
    \sigma^2\ge\frac{M_P}{6},
\]
the working-oracle errors are small enough for every ordinary trial
to satisfy the two acceptance tests. Interior trials with
$h=\epsilon$ also pass the high-accuracy validator.

\begin{lemma}[Safe scale]
\label{lem:gs-safe}
For fixed common constants $(\ell,\mu,\rho)$ and all sufficiently
small $\epsilon$, if
\[
    \sigma^2\ge\frac{M_P}{6},
\]
then every ordinary trial of
Algorithm~\ref{alg:gradient-scaled-scar} is accepted, and every
trial with
\[
    h=\epsilon,
    \qquad
    r<R(h,\sigma)
\]
passes validation.

At any generated scale $\sigma\ge\sigma_0$, successful validation
implies
\begin{equation}
    \|\nabla P(x^+)\|<\epsilon,
    \qquad
    \lambda_{\min}(\nabla^2P(x^+))
    \ge
    -\frac{85}{32}\sigma\sqrt\epsilon.
    \label{eq:gs-validator-return}
\end{equation}
\end{lemma}

The total-work argument uses a gradient-level potential. Define
\begin{equation}
    \mathcal L_\epsilon(h)
    :=
    \left\lceil
        \log_2(h/\epsilon)
    \right\rceil,
    \qquad
    \Phi
    :=
    \widehat P
    +
    \frac{\epsilon^{3/2}}{512\sigma}
    \mathcal L_\epsilon(h).
    \label{eq:gs-potential}
\end{equation}
For every ordinary accepted step, the analysis in
Appendix~\ref{app:persistent-scar-proofs} proves
\begin{equation}
    \Phi-\Phi^+
    \ge
    \frac{W}{32}
    +
    \frac{T}{1024}
    +
    \frac{U}{16}.
    \label{eq:gs-potential-decrease}
\end{equation}
Interior acceptance reduces the gradient level, whereas the
multiplier term $U$ pays for a possible increase after a boundary
step.

After accounting for the positive jumps caused by scale refreshes,
the same argument yields an explicit budget $B_\epsilon$ satisfying
\begin{equation}
    B_\epsilon
    =
    \Delta_P
    +
    O_{\rm fixed}
    \bigl(
        \epsilon^{3/2}\mathcal A_\epsilon^3
    \bigr)
    \label{eq:gs-budget-asymptotic}
\end{equation}
and
\begin{equation}
    \sum_{\rm accepted}
    \sigma_k^2\|d_k\|^3
    \le
    8B_\epsilon.
    \label{eq:gs-path-budget-main}
\end{equation}
This is the path budget that amortizes the variable tracking work.
Indeed, with
\[
    z_k:=\frac{\sigma_k\|d_k\|}{\sqrt\epsilon},
\]
the tracking rule satisfies
\[
    N_{{\rm tr},k}
    =
    \left\lceil
        \log_2(1+2^{14}z_k)
    \right\rceil
    \le
    16+z_k^3.
\]
Thus the accumulated tracking cost on accepted steps is controlled by
\eqref{eq:gs-path-budget-main}, rather than by multiplying the number
of outer iterations by an independent logarithmic inner cost.


\begin{theorem}[Persistent-SCAR parameter-free complexity]
\label{thm:gradient-scaled-scar}
Fix $\ell\ge\mu>0$ and $\rho>0$.
There exists
\[
    \epsilon_0
    =
    \epsilon_0(\ell,\mu,\rho)
    \in(0,1)
\]
such that the following holds simultaneously for every payoff
satisfying Assumption~\ref{ass:ncsc} with these common constants and
every initialization $(x_0,y_{-1})$.

Set, only for the purpose of stating the bound,
\[
    \kappa:=\frac{\ell}{\mu},
    \qquad
    M_P:=4\sqrt2\,\kappa^3\rho,
    \qquad
    \Delta_P:=P(x_0)-P^*.
\]
For every $0<\epsilon\le\epsilon_0$,
Algorithm~\ref{alg:gradient-scaled-scar} terminates while using only
$\epsilon$ and universal numerical constants.
Its output $x_{\rm out}$ satisfies
\begin{equation}
    \|\nabla P(x_{\rm out})\|<\epsilon,
    \qquad
    \lambda_{\min}(\nabla^2P(x_{\rm out}))
    \ge
    -\frac{85}{16\sqrt6}
    \sqrt{M_P\epsilon}.
    \label{eq:gs-return}
\end{equation}

If $Q$ denotes the total number of solved trust-region subproblems
and $N_y$ the total number of inner-gradient evaluations, including
the secant evaluations, then
\begin{subequations}
\label{eq:gs-complexities}
\begin{align}
    Q
    &=
    O\!\left(
        \Delta_P\sqrt{M_P}\,
        \epsilon^{-3/2}
    \right)
    +
    O_{\rm fixed}
    (\mathcal A_\epsilon^3),
    \label{eq:gs-outer-complexity}
    \\
    N_y
    &=
    O\!\left(
        \Delta_P\sqrt{\kappa M_P}\,
        \epsilon^{-3/2}
    \right)
    +
    O_{\rm fixed}
    (
        \sqrt\kappa\,
        \mathcal A_\epsilon^3
    ).
    \label{eq:gs-inner-complexity}
\end{align}
\end{subequations}
Here $O_{\rm fixed}$ may depend on the fixed payoff,
regularity constants, and initialization, but not on $\epsilon$.
Equivalently, the leading outer and inner terms are
\begin{equation}
    O\!\left(
        \Delta_P\kappa^{3/2}\sqrt\rho\,
        \epsilon^{-3/2}
    \right),
    \qquad
    O\!\left(
        \Delta_P\kappa^2\sqrt\rho\,
        \epsilon^{-3/2}
    \right).
    \label{eq:gs-leading-complexities}
\end{equation}

Using internal tolerance $(1536/7225)\epsilon$ gives the standard
$(\epsilon,\sqrt{M_P\epsilon})$ guarantee whenever the internal
tolerance lies in the stated regime.
\end{theorem}

The theorem removes the multiplicative target-accuracy logarithm from
the leading inner-gradient term. In contrast with
Corollary~\ref{cor:pf-scar-inner-complexity}, the persistent
curvature calibration is shared across the entire sequence of inner
objectives, and the variable residual reductions are charged to the
path budget \eqref{eq:gs-path-budget-main}.

\begin{remark}[Scope of the parameter-free guarantee]
\label{rem:pf-scope}
The threshold in
Theorem~\ref{thm:gradient-scaled-scar} is uniform over the class with
fixed common constants $(\ell,\mu,\rho)$ and over all
initializations, but it is not observable to the algorithm.
The theorem therefore does not provide an a posteriori
finite-accuracy certificate.
The distinction between this class-uniform statement and a guarantee
uniform over vanishing strong-concavity moduli is discussed in
Appendix~\ref{app:parameter-free-scope}.
\end{remark}

\begin{remark}[Oracle accounting]
\label{rem:pf-oracle-accounting}
The quantity $N_y$ counts evaluations of the inner gradient
$\nabla_yf$, which is the oracle counted in the SCAR complexity
result. SCAR's line search also evaluates the inner objective.
Under the standard implementation considered by
\citet{lan2026optimal}, these function-value evaluations have the
same leading order as the gradient count, up to an additive
line-search logarithm.

The corrected-gradient and Schur-complement constructions, including
the inverse actions involving $\nabla^2_{yy}f$, are not included in
$N_y$. If these inverse actions are implemented iteratively, their
accuracy and oracle complexity must be accounted for separately.
\end{remark}